\subsection{Warp Finance Case Study}
\label{subsec:Warp}


On Dec. 18, 2020, Warp Finance suffered a flash loan attack which leads to about \$7.8 million loss, 
according to Rekt~\cite{Rekt}. The first attack happened in 
the transaction~\cite{Warp}. 
A detailed postmortem analysis~\cite{WarpFiAnalysis} is written by SlowMist on Medium. 


\noindent \textbf{Attack Root Cause:}
The core of Warp attack is a design flaw of calculating 
the price of LP tokens. The price of LP token is calculated as (amount(WETH) in the pool * WETH price + amount(DAI) in the pool * DAI price) / total supply of LP. Even though the developer uses Uniswap official price oracle to (correctly) calculate the prices of WETH and DAI, they failed to realize the amounts of WETH and DAI can also be manipulated by flash loans. The exploiter took advantage of this point, pumped up the LP price, and falsely borrow excessive USDC and DAI from the pool.


\noindent \textbf{Real Attack Vector in History:}
\\
Action 1: \codeword{mint(2900030e18)} 
\\ 
Action 2: \codeword{swap(WETH, DAI, 341217e18)} 
\\ 
Action 3: \codeword{provideCollateral(94349e18)} 
\\ 
Action 4: \codeword{borrowSC(USDC, 3917983e6)} 
\\
Action 5: \codeword{borrowSC(DAI, 3862646e18)} 
\\
Action 6: \codeword{swap(DAI, WETH, 47622329e18)} 
\\ 
Adjusted profit: ~1,693,523 USD (DAI + USDC + WETH)


First the exploiter flash loaned WETH and DAI and \codeword{mint} 2900030 DAI and equivalent WETH to Uniswap's WETH-DAI pair. Then the attacker swapped a huge amount of WETH into DAI via the pair to increase the total value of the WETH-DAI pool, pumping up the unit price of LP token. The exploiter then mortgages the previously obtained LP Token through the \codeword{provideCollateral} function. As the unit price of LP Token becomes higher, the LP tokens mortgaged by the attacker give an abnormally high borrow limit thus the attacker can borrow more make profit.


\noindent \textbf{Best Attack Vector from {\name}:}
\\
Action 1: \codeword{swap(WETH, DAI, 480069e18)} 
\\
Action 2: \codeword{mint(322007e18)} 
\\
Action 3: \codeword{provideCollateral(65937e18)} 
\\ 
Action 4: \codeword{borrowSC(USDC, 3862530e6)} 
\\
Action 5: \codeword{borrowSC(DAI, 3917914e18)} 
\\
Action 6: \codeword{swap(DAI, WETH, 48862671e18)} 
\\ 
Adjusted profit: ~3,368,718 USD (DAI + USDC + WETH)


It is quite surprising that the attack vector from {\name} gets a higher profit compared to the attacker in history. The reason behind this is the attack vector given by {\name} spends more on manipulating Uniswap pair(Action 1 and Action 6) and spends less on the LP Token mortgaged(Action 2 and Action 3). In this attack, to get the same borrow limit(which can be used to borrow USDC and DAI in Action 4 and Action 5), it is more economical to spend on manipulating the Uniswap pair than collateralizing more LP tokens. 


